\chapter{Daily Routine, Stress and Prevention}
\label{ch:lifestyle-stress}

\emph{How does a yogic lifestyle organise the day and the year, how does yoga understand and manage stress, and what can it do to help prevent lifestyle disorders?}
This chapter covers three topics of the health unit that are about living well: the daily and seasonal routines of the Ayurvedic tradition, the yogic management of stress, and the place of yoga in the prevention of metabolic and respiratory disorders. Each is about prevention. None is about treatment.

\section{Dinacharya and ritucharya}

\subsection*{The idea}
The Ayurvedic view is that what is healthy depends on the person (age, constitution, condition), on the season and on the place. Daily and seasonal regimens keep the three doshas, vata, pitta and kapha, in balance and so prevent illness. A dosha moves through three stages: accumulation (\emph{sañcaya}), aggravation (\emph{prakopa}) and pacification (\emph{praśama}); the aim is balance, \emph{doṣa-sāmya}. Said simply: live in rhythm with the day and the season, and change gradually.

\subsection*{Dinacharya, the daily routine}
For a yogic life the day runs like this: wake early, in \emph{brāhma muhūrta}, roughly one to one and a half hours before sunrise (sources range from about 45 to about 96 minutes); clear the bowels and bladder; clean the teeth (\emph{danta pavana}), tongue (\emph{jihvā nirlekhana}), mouth and face (\emph{mukha prakṣālana}), eyes and nose; hold oil in the mouth or swish it (\emph{gaṇḍūṣa} or \emph{kavala}); exercise (\emph{vyāyāma}); massage with oil (\emph{abhyaṅga}); bathe (\emph{snāna}); practise yoga and meditation; take regular and moderate meals; work with a calm mind; and sleep at a regular time. The lists of items differ between texts.

The classical text divides each half of the day and of the night into three parts: kapha rules the first part, pitta the middle, vata the last. The clock hours often quoted (vata about 2 to 6, kapha about 6 to 10, pitta about 10 to 2) are a modern teaching, not in the classical text.

\subsection*{Ritucharya, the seasonal routine}
The year has six seasons in two halves. In \emph{ādāna kāla}, the taking-away half, strength tends to fall: \emph{śiśira} (winter), \emph{vasanta} (spring) and \emph{grīṣma} (summer). In \emph{visarga kāla}, the giving half: \emph{varṣā} (monsoon), \emph{śarad} (autumn) and \emph{hemanta} (late autumn or early winter). Diet and activity are adjusted to the season, so that the body is protected from its effect. For example, cold and dry weather raise vata, heat raises pitta, and cold, heavy, moist food raises kapha.

A commonly given table has kapha accumulating in shishira, aggravated in vasanta and settling in grishma; and pitta accumulating in varsha, aggravated in sharad and settling in hemanta.

\emph{Ṛtusandhi} is the junction of two seasons, about two weeks: the last seven days of one and the first seven of the next. Change diet and routine gradually then. Classical purification therapies are given only by a physician; for a personal routine, consult a qualified Ayurvedic practitioner. The Charaka Samhita says that one who knows and follows the seasonal regimen gains strength and a good complexion; the focus is prevention.

\section{Stress and its management}

\subsection*{What stress is}
Stress is the response of a person when the demands on them are greater than the resources to meet them. A mild, useful kind is \emph{eustress}; a harmful kind is \emph{distress}. Common causes are too many obligations and commitments, negative thinking and pessimism, and poor skill in handling conflict. The climate and culture of a workplace and the strain of dealing with other people, \emph{psychosocial} stress, are further sources. Frustration is the feeling of being upset because one cannot change or achieve something.

\subsection*{What stress does to the body}
Stress prepares the body for ``fight or flight'': breathing quickens, the heart beats faster, the blood pressure rises, blood is diverted to the muscles, the senses sharpen, the muscles tense, digestion slows, and stress hormones rise. This is the sympathetic nervous system at work (see \cref{ch:human-body}). Stress that is not resolved can end in psychosomatic disorders.

\subsection*{How the mind produces it}
Psychology is the scientific study of mental processes and behaviour. \emph{Overt} behaviour is what can be seen; \emph{covert} behaviour is the inner mental activity. The guidebook describes how a thought becomes a feeling: the senses (\emph{indriya}) bring in information, the intellect processes it with the help of memory, and emotion colours it. Thoughts are like waves on the sea.

The gunas shape the mind: sattva gives steadiness (love, compassion, honesty, calm); rajas gives activity (drive, ambition, desire); tamas gives inertia (laziness, sleep, dullness). Anger, fear, hatred and jealousy feed stress; peace, contentment and happiness refresh. In the yogic view the roots of stress are the five afflictions (\emph{kleśa}): ignorance, ego, liking, dislike and fear. Strong likes and dislikes upset the balance of the mind and give rise to the mental disturbance called adhi.

The Gita (second chapter, verses 62 and 63) traces the chain: dwelling on the objects of the senses leads to attachment, attachment to desire, desire to anger, and anger to confusion. Its rule for steady living is action without attachment, without liking or dislike, and without hunger for the fruit.

\subsection*{Mental health and mental hygiene}
Mental health, in the guidebook, is a state of well-being in which a person knows their abilities, copes with the usual stresses of life, works productively and contributes to the community. In Patanjali's terms yoga is the stilling of the vrittis by steady practice (\emph{abhyāsa}) and detachment (\emph{vairāgya}), through the inner limbs: pratyahara, dharana, dhyana and samadhi.

Mental hygiene is the habit of keeping the mind healthy by watching and managing one's thoughts and feelings. Its aims are to realise one's potential, to respect oneself and others, to accept the limits of oneself and others, to bring harmony and happiness, to adjust well to life, and to know oneself. Three yogic attitudes support it:
\begin{itemize}
\item \textbf{Pratipaksha bhavana}: cultivating the opposite thought to cancel a disturbing one.
\item \textbf{Anitya bhavana}: seeing that bodily experiences are impermanent, which grows detachment.
\item \textbf{Sakshi bhavana}: watching one's own actions as a witness, which grows self-awareness and an even mind.
\end{itemize}
The practices are said to widen awareness, with the traditional goal of union of the individual self with the supreme self; learn that as the yogic view.

\subsection*{What yoga offers}
The guidebook names four yogic means of managing stress: \emph{āhāra}, right food; \emph{vihāra}, proper recreation and rest; \emph{vicāra}, positive thinking; and \emph{vyavahāra}, correct conduct and action. In practice, yoga offers asanas and loosening practices that release tension; slow breathing and pranayama; meditation and relaxation; and right conduct and a regulated life.

Cyclic meditation, a stress-management technique based on the Mandukya Upanishad, alternates postures with rest lying on the back; the guidebook cites studies reporting better attention and emotional competence, lower oxygen consumption and lower anxiety. Om meditation makes the mind one-pointed and clears information overload. Prayer is said to bring relaxation, ease the strain of trying to control everything, raise positive feeling, and join body, mind and spirit. Karma yoga, action without attachment, removes the worry over results.

Say that yoga helps to manage stress; never promise the cure of a disorder.

\subsection*{Practical steps for a class}
Know your stress signs and your stressors. Breathe slowly. Practise regularly. Learn to let go and rest. Keep a balanced day.

\section{Prevention of metabolic and respiratory disorders}

\subsection*{The idea}
Metabolic and chronic respiratory disorders, such as obesity, diabetes, asthma and COPD, share risk factors such as tobacco, physical inactivity, unhealthy diet, harmful use of alcohol and air pollution (WHO). Yoga supports prevention as regular activity and a regulated, calm way of living. It offers gentle movement and flowing practices (surya namaskara and asanas), breath awareness and deep breathing, regulated diet and sleep, and relaxation and meditation. Say ``prevention and support'', never ``cure''.

\subsection*{The airways}
The respiratory system runs from the nose and throat to the lungs, driven by the diaphragm and the muscles that help breathing. Yoga values clear airways, and its cleansing practice for the nose is jala neti with warm salted water, a practice for well-being learned from a teacher.

\subsection*{Digestion and metabolism}
Good digestion is seen as the base of physical and mental well-being. The right food, in the right amount and at the right time, keeps it steady; eating for comfort, from greed or from emotional need upsets it. In yoga, the manipura chakra at the navel stands for the digestive organs and the inner fire (\emph{agni}) that digests. Asanas, pranayama and the cleansing actions are used to keep the abdomen healthy; dhauti, nauli and basti are the cleansing actions named for the digestive tract. They are learned only from a qualified teacher and are not for a beginner's home practice. The workshop's claim that these actions heal disorders is a treatment claim with no source and is left out.

\begin{caution}
Yoga is not a remedy for every illness and does not replace medical treatment. A person with a condition practises only after the doctor's advice and under a qualified teacher. Avoid kapalabhati and other forceful breathing in heart disease, high blood pressure, vertigo, epilepsy, migraine, stroke, hernia, gastric ulcer and nose bleeds; with any condition, practise only after the doctor's advice.
\end{caution}

\begin{recall}
\begin{itemize}
\item Dinacharya: daily routine; wake in brahma muhurta. Ritucharya: seasonal routine; ritusandhi is the seasonal junction.
\item Adana kala: shishira, vasanta, grishma. Visarga kala: varsha, sharad, hemanta.
\item Stress: demands greater than resources; eustress and distress.
\item Four yogic means: ahara, vihara, vichara, vyavahara.
\item Pratipaksha, anitya and sakshi bhavana.
\item Prevention and support, never cure.
\end{itemize}
\end{recall}

\subsection*{Worked examples}
\begin{itemize}
\item \textbf{An office worker sleeps late, skips breakfast and feels tense all day.} Apply dinacharya: a regular waking and sleeping time, regular moderate meals, a short morning practice of loosening, slow breathing and relaxation. Suggest a short break with slow breathing during the working day.
\item \textbf{A student keeps replaying an angry exchange with a colleague.} Pratipaksha bhavana: deliberately cultivate the opposite thought, such as goodwill or understanding; then a few minutes of slow abdominal breathing.
\item \textbf{A student with asthma asks to join kapalabhati.} Leave it out unless the doctor and the teacher agree; offer gentle deep breathing instead.
\end{itemize}

\subsection*{Check yourself}
\begin{enumerate}
\item What is dinacharya? What is ritucharya?
\item What is brahma muhurta?
\item Which seasons make up adana kala?
\item What is ritusandhi, and what is advised then?
\item Define stress in one line.
\item What are the two kinds of stress?
\item Name the four yogic means of stress management.
\item What is pratipaksha bhavana?
\item What is the difference between adhija and anadhija vyadhi?
\item Does yoga cure metabolic or respiratory disorders?
\item Which chakra stands for the digestive organs and the digestive fire?
\end{enumerate}

\subsection*{Answers}
\begin{enumerate}
\item The daily routine, lived in rhythm with the day; the seasonal routine, adjusting food and practice to the season.
\item The time about one to one and a half hours before sunrise, suggested for waking.
\item Shishira, vasanta and grishma.
\item The junction of two seasons, about two weeks; change diet and routine gradually.
\item The response of a person when the demands are greater than the resources to meet them.
\item Eustress (useful) and distress (harmful).
\item Ahara (right food), vihara (rest and recreation), vichara (positive thinking) and vyavahara (correct conduct).
\item Cultivating the opposite thought to cancel a disturbing one.
\item Adhija begins in the mind (the manomaya sheath); anadhija is a bodily ailment, mainly of the annamaya sheath.
\item No; it supports prevention and well-being and does not replace treatment.
\item Manipura.
\end{enumerate}

\paragraph{Sources.} The official guidebook for Level I (2016, older scheme), fifth and sixth chapters (stress, mental health, mental hygiene, cyclic meditation, the four means) and second chapter (the cleansing actions); the Bhagavad Gita, second chapter; Caraka Samhita online and published Ayurveda reviews on dinacharya and ritucharya (to be checked against a named Ayurvedic text); the Board syllabus for Level 1.
